\section{Block visits and epoch transcripts}
\label{sec:simulation}

Fix a public input \(x\) and cap \(T\ge2\) for which the access
condition in Section~\ref{sec:introduction} holds. All parameters below
are computed uniformly from \(T\). Let \(\rho\) be the fixed number of
read-only heads, and let \(\ell_\Sigma\ge1\) be a fixed number of bits
per writable-tape symbol. Use a fixed-width raw encoding of the
finite-control state, terminal flags, elapsed time, and all head
displacements. Its exact length can be chosen as
\[
 c=c_0+(\rho+2)\lceil\log_2(2T+3)\rceil=O(\log(T+2)),
\]
where \(c_0\) is a fixed constant for the state and flags. Unused field
encodings will be handled below. A halt at the cap is recorded as a
halt; otherwise reaching the cap is recorded as a timeout. After
either event, stationary dummy visits preserve the outcome.

\subsection{A visit accesses one block}

For an integer width \(b\ge1\) polynomially bounded in \(T\), and offset
\(a\in\{0,\ldots,b-1\}\), physical block \(v\in\mathbb Z\) is
\([a+vb,a+(v+1)b-1]\).
A visit runs until the writable head leaves its block, the machine
halts, or the cap is reached. The departure step belongs to the old
block: it writes before moving, and does not read the destination
cell until the next visit.

Put the writable origin at zero and write \(I_T=[-T,T]\cap\mathbb Z\).
The physical blocks meeting \(I_T\) have indices
\[
 v_-=\left\lfloor\frac{-T-a}{b}\right\rfloor,\qquad
 v_+=\left\lfloor\frac{T-a}{b}\right\rfloor.
\]
Relabel these intervals in increasing order as \(C_1,\ldots,C_N\), where
\begin{equation}\label{eq:block-count}
 N=v_+-v_-+1\le 2+\frac{2T}{b}.
\end{equation}
Signed division and translation between the two indexings use
\(O(\log(T+2))\) bits. The complete intervals \(C_i\) may extend
outside \(I_T\). Choose a fixed valid filler symbol \(s_{\mathrm{fill}}\)
with a canonical \(\ell_\Sigma\)-bit code. For a tape configuration
\(\Gamma\), its represented block \(\operatorname{rep}_{T,i}(\Gamma)\)
has the code of \(\Gamma(q)\) at \(q\in C_i\cap I_T\), and the filler
code at \(q\in C_i\setminus I_T\). This representation depends on the
trial cap; the underlying initial symbols remain cap-independent.

\begin{lemma}
\label{lem:visits}
Some offset has at most \(T/b\) block crossings in the capped
computation. For that offset at most
\(R=\lceil1+T/b\rceil\) visits reach the halt or timeout.
For every offset there are total maps on raw Boolean words
\[
 \operatorname{sel}_a:\bits^c\longrightarrow\{1,\ldots,N\},
 \qquad
 \operatorname{Visit}_a:
 \bits^c\times\bits^{b\ell_\Sigma}
 \longrightarrow \bits^c\times\bits^{b\ell_\Sigma}.
\]
On input \((z,B)\), \(B\) is interpreted as block
\(\operatorname{sel}_a(z)\). A visit changes only that block, agrees
with the capped computation on
represented actual data, and is computable in
\(O((b+1)\polylog T)\) work bits on explicit arguments.
\end{lemma}

\begin{proof}
Every unit head move crosses a boundary for exactly one offset;
a stationary move crosses none. The sum of crossing counts over all
\(b\) offsets is at most \(T\), proving the averaging claim.
The number of genuine visits is at most one plus the crossing count.

Fix a canonical encoding for each controller record. Accept a raw
controller only if its fields encode such a record, its elapsed time
satisfies \(0\le\tau\le T\), and every head displacement \(q_h\)
satisfies \(|q_h|\le\tau\). These checks preserve every actual
controller. On an accepted controller, \(\operatorname{sel}_a\)
uses the writable displacement and signed floor division to select
its label in \(\{1,\ldots,N\}\). On every rejected string it selects
the block containing the origin. Thus selection is total on raw
\(c\)-bit strings.

For a rejected controller, \(\operatorname{Visit}_a\) returns a fixed
canonically encoded terminal dummy controller with elapsed time and
all displacements zero,
selecting that designated block, and returns the raw block unchanged.
For an accepted terminal
controller it returns both arguments unchanged. In all other cases it
simulates the visit, recording halt before timeout if both occur at
the cap, and emits the outgoing controller in its fixed raw encoding.
At elapsed time \(T\) it executes no transition. The block is a buffer
of raw \(\ell_\Sigma\)-bit entries. Decode an entry only when the
simulated writable head accesses that cell, with a fixed symbol for
every invalid code; encode only the symbols written by simulated
transitions. All unvisited entries remain unchanged. These rules give
a fixed raw output even on malformed block arguments. The selected
interval \(C_i\) maps the writable position \(q\) to buffer entry
\(q-\min C_i\); the visit stops after a move leaves that interval.

The same space and access bounds hold for every accepted speculative
invocation. If a head begins at displacement \(q_0\) at elapsed time
\(\tau\), then after \(s\le T-\tau\) locally simulated steps,
\begin{equation}\label{eq:local-confinement}
 |q_s|\le |q_0|+s\le\tau+s\le T.
\end{equation}
A transition occurs only before the remaining budget is exhausted.
Consequently every simulated writable read and write stays in \(I_T\),
and every read-only query has \(|q_h|\le T\). The latter is answered by
\(\mathsf{Sym}(x,h,T,q_h)\); the full controller supplies every
read-only displacement. Rejected and terminal cases make no such
queries. In particular, no simulated transition in a local invocation
accesses a writable cell outside \(I_T\), so it preserves the exterior
filler. If \(h_{\mathrm{w}}\) is the writable symbol tag, the initial
represented block is generated by calling
\(\mathsf{Sym}(x,h_{\mathrm{w}},T,q)\) for its coordinates in \(I_T\)
and emitting the filler elsewhere. No larger-radius access condition
is used.

Induction on the simulated steps now shows that, from represented
actual data, the controller and all in-range symbols agree with the
capped computation, while the filler remains fixed. The buffer,
controller, and counters use \(O((b+1)\polylog T)\) bits, including
symbol access, and the remaining time budget bounds every local run.
This proves totality, agreement, and the claimed workspace.
\end{proof}

We try all offsets. After the allotted visits, an offset with neither
terminal flag is discarded. Every offset simulates the actual capped
trajectory, so any completed offset gives the same outcome.
At least one completes by Lemma~\ref{lem:visits}.
No transcript of all block addresses is stored.

\begin{lemma}[Local replay]\label{lem:local-replay}
Fix an integer \(0\le m\le b\). For a raw controller list
\(Z=(z_0,\ldots,z_m)\) and a raw block \(B\in\bits^{b\ell_\Sigma}\),
define
\[
 \operatorname{Replay}_i(Z,B)=(B^{\mathrm{out}},p_i)
\]
by the following bounded scan. Initialize a buffer to \(B\) and
\(p_i=1\). For \(t=1,\ldots,m\), if
\(\operatorname{sel}_a(z_{t-1})=i\), apply
\(\operatorname{Visit}_a\) to \(z_{t-1}\) and the buffer, retain its
new buffer, and conjoin to \(p_i\) the bitwise equality of its raw
\(c\)-bit output with \(z_t\). Continue the scan after a failed check.
Visits with a different selector leave this buffer and flag unchanged.
The buffer and flag are total Boolean functions of \(Z,B\), using
\(O((b+(m+1)c+1)\polylog T)\) local workspace and no word-evaluation
calls.

From any fixed raw \(z_0\) and raw blocks \(B_1,\ldots,B_N\), define a
reference evolution that applies the same total visit serially to its
selected block and leaves the other blocks unchanged. This global
\(N\)-block evolution specifies semantics only; the simulator never
stores or evaluates the tuple. For a proposed list with this \(z_0\),
all replay flags equal one if and only if the list is the unique raw
controller list of the reference evolution. When they equal one, each
returned buffer is the corresponding reference block after \(m\) visits.
\end{lemma}

\begin{proof}
The scan has a fixed number of visits, and every local visit is total
by Lemma~\ref{lem:visits}. It keeps one buffer, the explicit list and
short counters, giving the stated local bound. If the list is the
reference list, each replay applies exactly its block's reference
updates, so all checks pass and the buffers agree.

Conversely, suppose all flags equal one. Induct on \(t\), maintaining
that the proposed controllers through \(z_{t-1}\) and every replay
buffer after scanning the first \(t-1\) visits agree with the reference
evolution. This holds initially. The next incoming raw controller
therefore selects the same unique block \(i\). Its replay buffer is
the reference buffer, so the deterministic total visit produces the
reference outgoing raw controller and block. Its bitwise check forces
the proposed \(z_t\) to be that controller. Other buffers do not
change. This proves the induction and the claim for arbitrary raw
initial data.
\end{proof}

\subsection{Histories and checkpoints}
\label{sec:histories}

Set
\[
 k=\lceil T^{1/5}\rceil,\qquad b=k^2,\qquad
 E=\left\lceil\frac{\lceil1+T/b\rceil}{b}\right\rceil.
\]
For a fixed offset, take the first \(Eb\) visits, continuing with
stationary visits only after halt or timeout. There are
\(E=O(1+T/b^2)\) epochs, indexed by \(0\le e<E\), each consisting
of \(b\) visits, divided into
\(k\) rounds of \(k\) visits. An insufficient offset may still have
a nonterminal last controller.

Use the concrete common width
\begin{equation}\label{eq:word-width}
 d=\max\{1,b\ell_\Sigma,(b+1)c\}
   =O((b+1)\log(T+2))
\end{equation}
for the following Boolean words. Each type uses the indicated prefix
and pads the remaining bits with zero; on arbitrary words its parser
ignores those remaining bits.
\begin{itemize}
\item \(A_{e,i}\) is the padded represented block
\(\operatorname{rep}_{T,i}(\Gamma_e)\), where \(\Gamma_e\) is the actual
tape after the first \(eb\) visits of this offset.
\item \(H_{e,j}\), for \(0\le j\le k\), contains the epoch's incoming
raw controller and every outgoing raw controller of its first \(jk\)
visits.
\end{itemize}
The initial \(A_{0,i}\) and \(H_{0,0}\) are generated base words.
For \(e>0\), \(H_{e,0}\) copies the final raw controller slice of
\(H_{e-1,k}\) and adds the prescribed padding. For \(0\le e<E-1\),
parse \(H_{e,k}\) as its \(b+1\) raw slices and apply
\(\operatorname{Replay}_i\) to that list and the block prefix of
\(A_{e,i}\). Define \(A_{e+1,i}\) to be its buffer projection, padded
to width \(d\). Lemmas~\ref{lem:visits} and~\ref{lem:local-replay} prove
this checkpoint identity on semantic inputs. On malformed inputs the same projection
is still total, because the replay continues after failed checks.
These extraction and checkpoint rules have fan-in one and two,
respectively, and use \(O(d\polylog T)\) local workspace.
The target word is the final history \(H_{E-1,k}\).

\begin{lemma}
\label{lem:factored-history}
For \(0\le j<k\), each round update \(H_{e,j}\mapsto H_{e,j+1}\)
is a factored rule
of Lemma~\ref{lem:weighted-evaluation}, with distinguished child
\(H_{e,j}\), other children \(A_{e,1},\ldots,A_{e,N}\), and guess
length \(g=kc\le d\).
\end{lemma}

\begin{proof}
Write \(H=H_{e,j}\) and \(A_i=A_{e,i}\).
Parse the raw slices already in \(H\) as
\(z_0,\ldots,z_{jk}\), with the epoch-start controller appearing once
as \(z_0\). Write a candidate \(Y\in\bits^{kc}\) as \(k\) raw slices
\(y_1,\ldots,y_k\), and set \(z_{jk+s}=y_s\) for \(1\le s\le k\).
The resulting list
\[
 Z_Y(H)=(z_0,\ldots,z_{(j+1)k})
\]
specifies the entire epoch prefix through the new round. Define
\(G_Y(H)\) to copy this list into the padded history format, and
\(P_{Y,i}(H,A_i)\) to be the flag projection of
\(\operatorname{Replay}_i\) on \(Z_Y(H)\) and the block prefix of
\(A_i\). Thus replay starts at the epoch boundary, where \(A_i\) is
defined, and includes the preceding \(jk\) visits.

Fixed slicing, padding, and the total replay define total Boolean
functions even on malformed \(H,A_i,Y\). Since \((j+1)k\le b\),
Lemma~\ref{lem:local-replay} gives \(O(d\polylog T)\) local workspace,
including public-symbol access, and no word-evaluation calls.
At the semantic \(H,A_1,\ldots,A_N\), the fixed prefix agrees with
the reference evolution of the epoch, which agrees with the actual
capped computation by Lemma~\ref{lem:visits}. Lemma~\ref{lem:local-replay}
therefore says that all predicates pass for exactly one raw suffix
\(Y\), namely the actual next \(k\) outgoing controllers. Its padded
output is \(H_{e,j+1}\). Coordinatewise in characteristic two at
these arguments,
\[
 \overline{H_{e,j+1}}
 =\sum_{Y\in\bits^{kc}}\overline{G_Y(H)}
       \prod_{i=1}^N P_{Y,i}(H,A_i).
\]
The extra coordinate is one because exactly one passing indicator
is one.

An inconsistent fixed prefix may have no passing suffix; the evaluator
requires this identity only at semantic arguments. Its full homogeneous
lookup extensions are used for every predicate, including one that is
inactive at a particular Boolean history. This is the factored rule
of Lemma~\ref{lem:weighted-evaluation}.
\end{proof}

\subsection{The weighted dependency path}
\label{sec:path}

The word definitions are acyclic. Within an epoch put the start-array
words first, then \(H_{e,0},H_{e,1},\ldots,H_{e,k}\); all words of
epoch \(e-1\) precede those of epoch \(e\).
There are exactly \(E(N+k+1)\) word identifiers. Each consists of a
type and the indices \(e,i\) or \(e,j\), all of \(O(\log(T+2))\) bits.
The parent type and indices specify every child by binary arithmetic:
a checkpoint \(A_{e,i}\) has children \(A_{e-1,i},H_{e-1,k}\), an
extraction \(H_{e,0}\) has child \(H_{e-1,k}\), and a round word
\(H_{e,j}\), \(j\ge1\), has distinguished child \(H_{e,j-1}\) and
indexed children \(A_{e,1},\ldots,A_{e,N}\). Base words have no children.
This navigation uses \(\polylog T\) space without storing the graph.

The history-to-history edge has weight \(1\) in
Lemma~\ref{lem:weighted-evaluation}. A history-to-array edge has weight
\(g+1=kc+1\). Edges of the ordinary checkpoint and extraction rules
have weight \(1\). Figure~\ref{fig:epoch-path} displays the relevant
dependency structure.

\begin{figure}[htbp]
\centering
\begin{tikzpicture}[
  >=Latex,box/.style={draw,rounded corners=2pt,minimum width=29mm,minimum height=9mm},
  every node/.style={font=\small},
  cheap/.style={->,thick},
  costly/.style={->,very thick,dashed}
]
\node[box] (h) at (0,2) {$H_{e,j}$};
\node[box] (hp) at (5,2) {$H_{e,j-1}$};
\node[box] (a) at (0,0) {$A_{e,i}$};
\node[box] (ap) at (0,-2) {$A_{e-1,i}$};
\node[box] (oldh) at (5,-2) {$H_{e-1,k}$};
\draw[cheap] (h) -- node[above] {$1$} (hp);
\draw[costly] (h) -- node[left] {$kc+1$} (a);
\draw[cheap] (a) -- node[pos=0.3,left] {$1$} (ap);
\draw[cheap] (a) -- node[above,sloped] {$1$} (oldh);
\draw[dotted] (-2,-1) -- (8.7,-1);
\node[anchor=west] at (7,0) {epoch \(e\)};
\node[anchor=west] at (7,-2) {epoch \(e-1\)};
\end{tikzpicture}

\caption{Representative dependencies for \(e\ge1\) and
\(1\le j\le k\). A path may follow preceding histories before entering
\(A_{e,i}\); the omitted extraction edge
\(H_{e,0}\to H_{e-1,k}\) is another route to the preceding epoch.
The weight \(kc+1\) charges the round guess retained in a suspended
frame. Every child of that array word belongs to the preceding epoch,
so the same path cannot retain another round guess in epoch \(e\).
All weights bound suspended storage up to a common polylogarithmic factor.}
\label{fig:epoch-path}
\end{figure}

\begin{lemma}
\label{lem:epoch-weight}
Every dependency path from the final history has weight
\[
 O(E(k+kc+1))=O(Ek\log(T+2)).
\]
\end{lemma}

\begin{proof}
In epoch \(e\), a path takes at most \(k\) preceding-history edges.
It may then use the extraction edge from \(H_{e,0}\), or enter a
start-array word \(A_{e,i}\) and pay \(kc+1\).
For \(e>0\), that array word has children only \(A_{e-1,i}\) and
\(H_{e-1,k}\); for \(e=0\), it is a base word.
Thus a path has at most one expensive edge and a constant number of
ordinary edges in each epoch. Paths entering an epoch at an array
word satisfy the same bound. Summing over the epochs proves the
claim. Sequential repetitions and reverse calls follow the same
dependencies and do not increase the maximum active path weight.
\end{proof}

\begin{proof}[Proof of Theorem~\ref{thm:main}]
Apply Lemma~\ref{lem:weighted-evaluation} to the base, checkpoint,
extraction and factored history words.
The concrete width in \eqref{eq:word-width} covers each block and
history, and \(g=kc\le d\) because \(b=k^2\ge k\). The word count
\(E(N+k+1)\) is polynomial in \(T\) by \eqref{eq:block-count} and the
chosen parameters. Ordinary fan-in is at most two, so the evaluator
uses three vector registers. For this offset one may compute and supply
the degree bound
\[
 \Delta=(N+1)d
 =O((T+b+1)\log(T+2))=T^{O(1)},
\]
which covers the ordinary degree \(d\) and every history degree.
The candidate count \(2^{kc}\) and the Boolean lookup enumerations
affect time, not this degree or the number of word identifiers.
Base coordinates are fixed controller bits or use the accessor or
filler. Extraction and appending copy fixed slices, and each replay
uses one block and at most \(b+1\) controller slices. The local
Boolean kernels therefore use
\(O(d\polylog T)\) local space and make no word-evaluation calls.
The preceding navigation description supplies the required uniformity.
\mbox{Lemmas~\ref{lem:factored-history} and~\ref{lem:epoch-weight}} verify
the remaining hypotheses. Starting with a zero target yields
the final history, hence its final controller, in space
\[
 \softO(d+Ek)=
 \softO\!\left(b+(1+T/b^2)k\right)
 =\softO(T^{2/5}).
\]
The register pool, reusable local workspace and suspended stack are
all counted. In particular, the current unsuspended guess is included
in the word-sized local term, and no array of checkpoints is materialized.
Trying the offsets reuses this storage, and a completed offset
exists. Reading its finite-control state and halt/timeout flags gives
the required capped outcome.

For the unknown-running-time clause, the access condition holds for
\((x,T)\) at every \(T\ge2\), so it applies separately to all trials
with caps \(2,4,8,\ldots\). Continue after a timeout and return only on
an actual halt. If the machine halts in \(t\) steps, the first sufficient
cap is at most \(2\max(t,2)\). Each trial restarts from the same
cap-independent initial symbols, using its own represented blocks, and
discards its workspace before the next. This gives the asserted bound.
\end{proof}

\begin{remark}
\label{rem:output}
A specified bit of a sequential write-only output can also be
computed: keep its index, an output counter and the monitored bit in
the controller. For indices in the capped output range this costs
\(O(\log(T+2))\) bits. Enlarge \(c\) by these fields and use the same
width formula \eqref{eq:word-width}; the estimates are unchanged.
This does not require storing the output stream or change which
writable block a visit accesses.
\end{remark}

\subsection{What varying the three scales achieves}
\label{sec:balance}

Let the block width be \(B\), epoch length in visits be \(L\), and
round length be \(q\). Ignoring rounding and logarithms, a block
contributes \(B\) bits and the complete epoch transcript contributes
\(L\) bits to the materialized word size \(B+L\). There are
\(T/(BL)\) epochs and \(L/q\) rounds per epoch. The preceding-history
edges therefore contribute \(T/(Bq)\) to the weighted stack, while
one suspended \(q\)-controller guess per epoch contributes
\(Tq/(BL)\). Thus the stack estimate is
\[
 \frac{T}{BL}\left(\frac Lq+q\right)
 =\frac{T}{Bq}+\frac{Tq}{BL}.
\]
At \(q\simeq\sqrt L\), this becomes \(2T/(B\sqrt L)\).
If both \(B,L\le S\), it is at least \(2T/S^{3/2}\).
Balancing this displayed estimate with storage \(S\) therefore gives
the two-fifths exponent. Boundary and ceiling terms are nonnegative
and do not improve this balance. This calculation concerns this
estimate only: it rules out neither improved accounting nor another
transcript or simulation method.
